\color{unimain}{\subsection{Контроль положения обходного выключателя}}
\color{black}

\begin{enumerate}

\item
Контроль положения обходного выключателя ВН вводится накладкой \textbf{«КПОВ ВН»}.
Контроль положения обходного выключателя СН вводится накладкой \textbf{«КПОВ СН»}.

\begin{figure}[H]
\centering
\includegraphics[width=0.6\textwidth,height=0.6\textheight,keepaspectratio]{\fbpath/762.3481 КПОВ/img/KPOV.pdf}
\captionof{figure}{Функциональный блок КПОВ}\label{pic:KPOV}
\end{figure}

\small
\begin{longtable}{|>{\centering\arraybackslash}m{5.3cm}|>{\centering\arraybackslash}m{4.2cm}|>{\centering\arraybackslash}m{1cm}|>{\centering\arraybackslash}m{3.3cm}|>{\centering\arraybackslash}m{1.8cm}|}
% Первая страница: полный заголовок
\caption{Перечень уставочных параметров КПОВ\hfill\vspace{-0.5\baselineskip}}\label{desc:tbl1}\\
\hline
\rowcolor{gray!30}
Наименование уставки & Диапазон & Шаг & Значение по умолчанию & Единица измерения \\
\hline
\endfirsthead

% Промежуточные страницы: «Продолжение таблицы»
\caption*{\hspace{3pt}\emph{Продолжение таблицы \ref{desc:tbl1}}\hfill\ }\vspace{-0.5\baselineskip} \\
\hline
\rowcolor{gray!30}
Наименование уставки & Диапазон & Шаг & Значение по умолчанию & Единица измерения \\
\hline
\endhead

%===m>ДЗАТ|-
%===m
\end{longtable}
\normalsize	

\end{enumerate}